\section{Discussion}
\label{sec:discussion}

\paragraph{What to report.}
A multi-seed mean alone does not answer whether a federated LoRA ranking is stable under the non-IID draw. FFA-LoRA already shows that across-run standard deviations can exceed method gaps~\cite{sun2024ffalora}; our results go one step further by decomposing that variation into a partition component, a partition-by-method interaction, and a residual training-seed component. At Dirichlet $\alpha{=}0.1$, the residual $\hat\sigma^2_E$ is small relative to $\hat\sigma^2_P$ and $\hat\sigma^2_{PM}$ on TinyLlama, so reporting only run seeds understates the uncertainty that matters for ranking. Practical reporting should therefore (i) separate the partition seed from the training seed, (ii) publish the variance shares or at least the paired per-draw SD of method differences, and (iii) state how many partition draws were used. When the observed gap is near that paired SD ($0.00282$ on TinyLlama, $0.00473$ on LLaMA-3.2-3B), a few-draw evaluation is nearly a coin flip; when the gap is about $0.023$ to $0.031$, three paired draws sufficed here.

\paragraph{What pairing fixes, and what it does not.}
Pairing methods on the same partition draws removes the shared partition main effect from the comparison and is the right default for ranking. It does not remove the interaction: methods can still reorder across draws when their responses to a partition differ. That is exactly the FedIT--FLoRA near-tie ($+0.00055$ on TinyLlama, $-0.0020$ on LLaMA), where best-method and full-order flip probabilities coincide because only the top pair moves. Pairing also does not invent power. Detecting a pre-registered $\Delta{=}0.005$ still needs $5$ (TinyLlama) or $10$ (LLaMA) paired draws; the observed near-tie would need $206$ and $46$. Unpaired designs pay a similar or larger sample-size bill. Pairing is therefore a variance-control device, not a substitute for enough independent partitions.

\paragraph{Cross-scale reading.}
On TinyLlama the partition share is substantial ($0.508$). On LLaMA-3.2-3B the interaction dominates ($0.764$) and the partition main-effect CI includes zero, so we do not claim a partition main effect at 3B. Cross-scale CIs on the partition share overlap; the two scales are not distinguishable on that margin. The rise in LLaMA flip probability from $k{=}1$ to $k{=}3$ is a six-partition near-tie artifact, not evidence that more draws hurt. We report the 3B IID floor descriptively only ($2$ degrees of freedom per method) and do not interpret its ratio to $\hat\sigma^2_E$.

\paragraph{Limitations.}
The study uses one dataset (Dolly), one heterogeneity level ($\alpha{=}0.1$), three aggregators, and held-out next-token loss rather than generation benchmarks. LLaMA production holdouts run in float16; the float32 spot-check bound the difference well below the $0.01$ effect scale, but the primary 3B numbers remain float16. Only $6$ partitions support the 3B grid, which widens bootstrap CIs and limits ranking resamples. The MixedLM partition component did not identify, so the mixed-model cross-check is partial and moment estimates are primary. Communication is reported within method only, because FFA-LoRA's implemented payload differs by construction. RQ4 correlations are exploratory and do not replicate at 3B. Adapter weights were not retained for most cells, so the release covers results and metadata rather than full artifacts.

\paragraph{No method recommendation.}
We do not crown FedIT, FFA-LoRA, or FLoRA. The separable FFA-LoRA gaps are large and stable; the FedIT--FLoRA gap is not. The contribution is the measurement design and the conditioning of ranking stability on effect size relative to partition-draw noise.
